\chapter*{Заключение}                       % Заголовок
\addcontentsline{toc}{chapter}{Заключение}  % Добавляем его в оглавление

%% Согласно ГОСТ Р 7.0.11-2011:
%% 5.3.3 В заключении диссертации излагают итоги выполненного исследования, рекомендации, перспективы дальнейшей разработки темы.
%% 9.2.3 В заключении автореферата диссертации излагают итоги данного исследования, рекомендации и перспективы дальнейшей разработки темы.
%% Поэтому имеет смысл сделать эту часть общей и загрузить из одного файла в автореферат и в диссертацию:

% Основные результаты диссертационной работы заключаются в следующем:
% \input{common/concl}

В диссертационной работе решена актуальная научно-техническая задача разработки метода поиска и навигации по цифровым архивам древнекитайских рукописных документов на основе структурного представления иероглифических изображений. Актуальность исследования определяется высокой вариативностью древних рукописей, наличием большого числа редких иероглифов, ограниченным объёмом размеченных данных и недостаточной эффективностью существующих методов автоматической расшифровки при работе с историческими архивами.

В ходе исследования выполнен систематический анализ современных методов обработки рукописных документов и показано, что существующие подходы либо требуют больших объёмов обучающих данных, либо обладают недостаточной устойчивостью к особенностям древних рукописей. На основании проведённого анализа обоснована постановка задачи контекстного поиска по изображению иероглифа, в которой основное внимание уделяется поиску и навигации по архивным документам, а не полной расшифровке текста.

Для решения поставленной задачи разработан новый метод структурного представления древнекитайских иероглифов, основанный на использовании непрерывной морфологии и медиального представления формы. Предложен алгоритм построения скелетного метаграфа и формирования компактного множества критических точек, интерпретируемого как «созвездие» признаков. Полученное описание сохраняет основные геометрические и топологические свойства символов, обладает устойчивостью к изменениям масштаба, индивидуальным особенностям почерка и локальным дефектам исторических документов, а также обеспечивает компактное и интерпретируемое представление структуры иероглифов.

На основе разработанного признакового описания предложен новый метод вычисления меры сходства между изображениями, основанный на решении задачи о назначениях средствами линейного программирования. Разработаны две взаимодополняющие оптимизационные модели сопоставления критических точек, дополненные ограничениями, обеспечивающими топологическую согласованность соответствий. Предложенный подход позволяет эффективно сравнивать изображения иероглифов без предварительного обучения и сохраняет возможность обработки ранее неизвестных классов символов, что особенно важно для древнекитайских рукописных архивов.

Практическая эффективность разработанного метода подтверждена экспериментальными исследованиями на наборах данных CASIA-AHCDB и MTHv2. В задаче однократного распознавания предложенный подход достиг точности 97,23\% для сценария «5-way» и 95,84\% для «20-way», продемонстрировав результаты, сопоставимые и в ряде случаев превосходящие современные нейросетевые методы. При решении задачи поиска по архивным документам получены значения mAP = 0,7944, Recall@10 = 0,7455, Recall@50 = 0,9, Recall@100 = 0,93 и Precision@10 = 0,77, что свидетельствует о высокой эффективности ранжирования и успешном обнаружении большинства релевантных изображений в верхней части списка результатов поиска. Кроме того, показано, что предложенный алгоритм не требует предварительного обучения и может эффективно выполняться на стандартном персональном компьютере без использования специализированных графических ускорителей.

Таким образом, поставленная цель исследования достигнута, а все сформулированные задачи успешно решены. Научная новизна работы заключается в разработке нового метода структурного представления древнекитайских иероглифов на основе непрерывной морфологии, создании оптимизационных моделей вычисления меры сходства между такими представлениями и построении поисково-ориентированного подхода, предназначенного для работы с открытым множеством символов в условиях ограниченного объёма обучающих данных.

Перспективы дальнейших исследований связаны с расширением набора используемых структурных признаков, интеграцией разработанного подхода с современными методами представления изображений, а также распространением предложенной технологии на обработку других видов исторических рукописных документов.

В заключение автор выражает благодарность и большую признательность научному руководителю доктору технических наук Местецкому Л.М. за поддержку, помощь, обсуждение результатов и научное руководство. Автор также благодарит много разных людей и~всех, кто сделал настоящую работу автора возможной.
